\section{JML e verifica statica (ESC)}
\label{sec:jml-ref}

Per la parte di design by contract ho annotato con JML una classe sola,
\texttt{Parcheggio}. È il nucleo del dominio: tiene i due contatori dei posti liberi e
gestisce l'assegnazione e il rilascio, cioè fa quello che nel modello ASM fanno le regole
\texttt{r\_gestione\_VER}, \texttt{r\_gestione\_INGR} e \texttt{r\_gestione\_USC}.

L'ho tenuta piccola apposta e senza dipendenze da niente: dentro ci sono solo
\texttt{int} ed enum. Il motivo è pratico, cioè che così OpenJML riesce a dimostrarla
tutta, mentre su una classe che si tira dietro mezzo Spring non ci sarei mai arrivato.
Tutto il resto del sistema (la FSM, la UI, il display, i sensori) è Java normale senza
annotazioni, come chiedeva la consegna.

Una cosa che ho scoperto sbattendoci contro: i campi \texttt{postiStd} e
\texttt{postiDis} li ho dovuti dichiarare \texttt{private /*@ spec\_public @*/}. Le
clausole di specifica public, come le invarianti e le postcondizioni, possono nominare
solo roba visibile a livello public, e se lascio i campi semplicemente private OpenJML si
ferma con un errore di visibilità. Con \texttt{spec\_public} restano private per il
codice ma diventano visibili alla specifica.

\subsection{Invarianti}

Le invarianti sono due e devono valere sempre, da dopo il costruttore in poi:

\begin{itemize}
\item i posti standard liberi stanno sempre tra 0 e \texttt{MAX\_STD};
\item i posti disabili liberi stanno sempre tra 0 e \texttt{MAX\_DIS}.
\end{itemize}

\begin{lstlisting}[style=terminal, caption={Invarianti di classe}]
//@ public invariant 0 <= postiStd && postiStd <= MAX_STD;
private /*@ spec_public @*/ int postiStd;

//@ public invariant 0 <= postiDis && postiDis <= MAX_DIS;
private /*@ spec_public @*/ int postiDis;
\end{lstlisting}

Sono le stesse identiche proprietà che avevo già verificato con il model checking sul
modello, cioè \texttt{ag(posti >= 0)} e \texttt{ag(posti <= 1)}. Le ho scritte apposta
uguali: così la stessa cosa finisce garantita a tre livelli diversi, nel modello con il
model checker, nel codice con il contratto, e ancora nei test con JaCoCo.

\subsection{Costruttori}

I costruttori sono due. Quello di default garantisce lo stato iniziale del modello ASM,
cioè tutti i posti liberi. Quello con i parametri lo uso nei test, quando mi serve
partire da una configurazione qualunque, e chiede che i valori passati siano già dentro i
limiti dell'invariante.

\begin{lstlisting}[style=terminal, caption={Contratti dei costruttori}]
//@ ensures postiStd == MAX_STD && postiDis == MAX_DIS;
public Parcheggio() { ... }

//@ requires 0 <= postiStdIniziali && postiStdIniziali <= MAX_STD;
//@ requires 0 <= postiDisIniziali && postiDisIniziali <= MAX_DIS;
//@ ensures postiStd == postiStdIniziali && postiDis == postiDisIniziali;
public Parcheggio(int postiStdIniziali, int postiDisIniziali) { ... }
\end{lstlisting}

Il \texttt{requires} del secondo costruttore dice cosa deve garantire chi chiama, ma non
protegge a runtime: se qualcuno passa un valore fuori range senza aver fatto la verifica
statica, il contratto non lo ferma. Per questo dentro ho messo anche un controllo
esplicito che lancia \texttt{IllegalArgumentException}, e quel comportamento lo provo in
\texttt{ParcheggioTest} con \texttt{assertThrows}.

\subsection{Metodo \texttt{assegnaPosto}}

Questo è quello con il contratto più lungo, perché deve dire tutto quello che fa
\texttt{r\_gestione\_VER}. Le postcondizioni sono due, una per ramo.

Per STD e ABBONATO ci sono due esiti possibili: o c'era un posto standard libero, e
allora torna \texttt{POSTO\_STD} con il contatore sceso di 1, oppure non c'era e torna
\texttt{NESSUNO} senza toccare niente.

Per il DISABILE gli esiti sono tre, perché in mezzo c'è il fallback: posto disabile se
disponibile, altrimenti posto standard se disponibile, altrimenti \texttt{NESSUNO}.

\begin{lstlisting}[style=terminal, caption={Contratto di assegnaPosto}]
//@ requires utente != null;
//@ ensures (utente == TipoUtente.STD || utente == TipoUtente.ABBONATO) ==>
//@     ((\old(postiStd) > 0 && \result == TipoPosto.POSTO_STD
//@       && postiStd == \old(postiStd) - 1 && postiDis == \old(postiDis))
//@      || (\old(postiStd) == 0 && \result == TipoPosto.NESSUNO
//@       && postiStd == \old(postiStd) && postiDis == \old(postiDis)));
//@ ensures utente == TipoUtente.DISABILE ==>
//@     ((\old(postiDis) > 0 && \result == TipoPosto.POSTO_DIS
//@       && postiDis == \old(postiDis) - 1 && postiStd == \old(postiStd))
//@      || (\old(postiDis) == 0 && \old(postiStd) > 0
//@       && \result == TipoPosto.POSTO_STD
//@       && postiStd == \old(postiStd) - 1 && postiDis == \old(postiDis))
//@      || (\old(postiDis) == 0 && \old(postiStd) == 0
//@       && \result == TipoPosto.NESSUNO
//@       && postiStd == \old(postiStd) && postiDis == \old(postiDis)));
public TipoPosto assegnaPosto(TipoUtente utente) { ... }
\end{lstlisting}

Il \texttt{\textbackslash old} serve per parlare del valore che il contatore aveva prima
della chiamata, altrimenti non potrei dire ``è sceso di uno''.

Una cosa che all'inizio non avevo messo: in ogni esito ho anche scritto che l'\emph{altro}
contatore resta com'era. Senza quel pezzo il contratto sarebbe stato sottospecificato,
cioè un'implementazione che assegna il posto giusto ma nel frattempo tocca anche l'altro
contatore lo rispetterebbe lo stesso.

\subsection{Metodo \texttt{liberaPosto}}

Questo corrisponde al rilascio del posto in \texttt{r\_gestione\_USC}. Chiede che il
posto non sia null e che il contatore corrispondente non sia già al massimo, e garantisce
che venga incrementato di 1 solo quello giusto. Con \texttt{NESSUNO} non fa niente e lo
stato resta identico.

\begin{lstlisting}[style=terminal, caption={Contratto di liberaPosto}]
//@ requires posto != null;
//@ requires posto == TipoPosto.POSTO_STD ==> postiStd < MAX_STD;
//@ requires posto == TipoPosto.POSTO_DIS ==> postiDis < MAX_DIS;
//@ ensures posto == TipoPosto.POSTO_STD ==>
//@     postiStd == \old(postiStd) + 1 && postiDis == \old(postiDis);
//@ ensures posto == TipoPosto.POSTO_DIS ==>
//@     postiDis == \old(postiDis) + 1 && postiStd == \old(postiStd);
//@ ensures posto == TipoPosto.NESSUNO ==>
//@     postiStd == \old(postiStd) && postiDis == \old(postiDis);
public void liberaPosto(TipoPosto posto) { ... }
\end{lstlisting}

Quel secondo e terzo \texttt{requires} non c'erano nella prima versione e sono il motivo
per cui la prova non chiudeva: ne parlo tra poco.

\subsection{Metodi getter}

I due getter li ho annotati \texttt{/*@ pure @*/}, che vuol dire senza effetti
collaterali. Serve perché un metodo non puro dentro una clausola JML non si può usare. La
postcondizione dice solo che restituiscono il campo, niente di più.

\begin{lstlisting}[style=terminal, caption={Getter puri}]
//@ ensures \result == postiStd;
public /*@ pure @*/ int getPostiStd() { ... }

//@ ensures \result == postiDis;
public /*@ pure @*/ int getPostiDis() { ... }
\end{lstlisting}

\subsection{Verifica dei contratti con OpenJML (ESC)}

Scritti i contratti, li ho verificati con OpenJML (release 21.0.27, build nativa per
macOS arm64) in modalità \texttt{--esc}. Quello che fa il tool è tradurre codice e
contratti in formule logiche e passarle al solver Z3, che deve dimostrare che
\emph{qualunque} esecuzione rispetta postcondizioni e invarianti. È una cosa diversa dai
test: i test provano alcuni casi, qui invece o si dimostra per tutti o non si dimostra.

Il comando, dalla root del progetto, è questo:

\begin{lstlisting}[style=terminal, caption={Comando di verifica ESC}]
$ openjml --esc --progress \
    -cp java/src/main/java \
    java/src/main/java/it/tvsw/smartparking/core/Parcheggio.java
\end{lstlisting}

\subsubsection*{Un contratto che all'inizio non chiudeva}

Arrivare a ``tutti Verified'' non è stato immediato, e il caso di \texttt{liberaPosto}
secondo me è quello che spiega meglio come funziona la cosa.

La prima versione che avevo scritto era questa, con la postcondizione che descrive
l'incremento ma senza il \texttt{requires} sulla capacità:

\begin{lstlisting}[style=terminal, caption={Prima versione di liberaPosto (senza il requires di capacità)}]
//@ requires posto != null;
//@ ensures posto == TipoPosto.POSTO_STD ==>
//@     postiStd == \old(postiStd) + 1 && postiDis == \old(postiDis);
public void liberaPosto(TipoPosto posto) { ... }
\end{lstlisting}

Il solver, non avendo nessun vincolo sul valore di partenza, prova anche il caso
\texttt{postiStd == MAX\_STD}. Lì \texttt{postiStd + 1} fa \texttt{MAX\_STD + 1}, che
viola l'invariante di classe. Quindi non riesce a ristabilire l'invariante all'uscita del
metodo e mi risponde così:

\begin{lstlisting}[style=terminal, caption={Errore ESC sulla prima versione}]
Parcheggio.java:.. warning: The prover cannot establish an assertion
   (InvariantExit) in method liberaPosto
   postiStd == \old(postiStd) + 1 ...
   Associated declaration: invariant 0 <= postiStd && postiStd <= MAX_STD
\end{lstlisting}

La soluzione è rafforzare la precondizione, in modo che sia chi chiama a garantire che
c'è spazio per incrementare:

\begin{lstlisting}[style=terminal, caption={Versione corretta: il requires di capacità chiude la prova}]
//@ requires posto == TipoPosto.POSTO_STD ==> postiStd < MAX_STD;
//@ requires posto == TipoPosto.POSTO_DIS ==> postiDis < MAX_DIS;
\end{lstlisting}

Con \texttt{postiStd < MAX\_STD} assunto all'ingresso, dopo l'incremento vale
\texttt{postiStd <= MAX\_STD}, l'invariante regge e il metodo passa a Verified.

Il senso della cosa, secondo me, è questo: ogni volta che c'è un'operazione aritmetica
bisogna aver vincolato prima tutto quello che ci entra, con un \texttt{requires} o con
un'invariante. \texttt{Parcheggio} chiude tutto perché ha due soli interi e tutti e due
sono limitati; su un nucleo con array o contatori liberi la stessa cosa diventa molto più
difficile. L'altro ostacolo che ho incontrato era di tutt'altro tipo, cioè la visibilità
dei campi private, e l'ho risolto con \texttt{spec\_public} come ho detto all'inizio.

\screenshot{openjml_esc_output}{Output di \texttt{openjml --esc}: tutti i metodi Verified}{%
Esito Verified per tutti e sei i metodi di Parcheggio, senza righe ``cannot be proved''.}

Alla fine tutti e sei i metodi (i due costruttori, \texttt{assegnaPosto},
\texttt{liberaPosto} e i due getter) risultano Verified.
